% TikZ 图：环形电流与通电螺线管的磁场方向
% Day 21：磁场与安培力
\begin{tikzpicture}[scale=0.88,>=Latex,line cap=round,line join=round]
  \path[use as bounding box] (-5.10,-2.05) rectangle (5.10,2.05);

  % 左图：正视环形电流，逆时针电流在中心产生向外的磁场。
  \begin{scope}[shift={(-2.65,0)}]
    \node[font=\small] at (0,1.70) {环形电流：此面为 $N$ 极};
    \draw[red,very thick] (0,0) circle (1.12);
    \draw[red,very thick,->] (1.12,0) arc (0:72:1.12);
    \draw[red,very thick,->] (-1.12,0) arc (180:252:1.12);
    \node[red,font=\small] at (0,-1.45) {$I$ 沿逆时针方向};

    \draw[blue,thick] (0,0) circle (0.18);
    \fill[blue] (0,0) circle (0.052);
    \node[blue,font=\small] at (0,0.42) {$B$ 向外};
  \end{scope}

  % 右图：从左端看电流逆时针，因此左端为 N 极，管内磁场向左。
  \begin{scope}[shift={(2.55,0)}]
    \node[font=\small] at (0,1.70) {通电螺线管};
    \draw[thick] (-1.65,-0.52) rectangle (1.65,0.52);
    \foreach \x in {-1.45,-1.10,-0.75,-0.40,-0.05,0.30,0.65,1.00,1.35} {
      \draw[red,thick] (\x,-0.68)
        .. controls (\x+0.20,-0.24) and (\x+0.20,0.24) .. (\x,0.68);
    }

    \node[font=\small] at (-1.93,0) {$N$};
    \node[font=\small] at (1.93,0) {$S$};
    \draw[blue,very thick,->] (1.35,0) -- (-1.35,0);
    \node[blue,font=\small] at (0,0.30) {管内 $B$};

    \draw[blue,thick,->] (-1.55,0.82)
      .. controls (-0.45,1.35) and (0.45,1.35) .. (1.55,0.82);
    \draw[blue,thick,->] (-1.55,-0.82)
      .. controls (-0.45,-1.35) and (0.45,-1.35) .. (1.55,-0.82);

    \node[red,font=\scriptsize,align=center] at (0,-1.58)
      {从左端看电流逆时针\\四指沿电流，拇指指向左端};
  \end{scope}
\end{tikzpicture}
